
\subsubsection{3.1 Hypothesis Framework}

\textbf{H-E1 (Existence):} Class-wise accuracy profiles exhibit variance beyond what overall accuracy and per-class difficulty explain.

\textbf{H-M1 (Mechanism):} Weight matrix features correlate with class-wise accuracy profiles, explaining variance beyond the stratified baseline.

\subsubsection{3.2 Behavioral Variance Decomposition (H-E1)}

A stratified baseline predicts class-wise accuracy from overall model accuracy:

$$\hat{a}_{i,c} = A_i \cdot \frac{d_c}{\bar{d}}$$

where $A_i$ is overall accuracy of model $i, d_c$ is mean accuracy on class $c$ across all models, and $\bar{d}$ is mean of class difficulties.

The residual variance ratio measures the fraction not explained:

$$\text{Residual Ratio} = \frac{\text{Var}(a_{i,c} - \hat{a}_{i,c})}{\text{Var}(a_{i,c})}$$

A ratio > 0.05 indicates meaningful behavioral variance beyond what the stratified baseline captures.

\subsubsection{3.3 Weight Statistics Probe (H-M1)}

For each model, 5 statistics (mean, standard deviation, minimum, maximum, L2 norm) are extracted from 5 layers (3 convolutional, 2 fully-connected), yielding 25 features. The layer keys are:
\begin{itemize}
\item \texttt{module\textbackslash \_list.0.weight} (conv1)
\item \texttt{module\textbackslash \_list.3.weight} (conv2)
\item \texttt{module\textbackslash \_list.6.weight} (conv3)
\item \texttt{module\textbackslash \_list.9.weight} (fc1)
\item \texttt{module\textbackslash \_list.11.weight} (fc2)
\end{itemize}

Ridge regression ($\alpha$ = 1.0) predicts per-class accuracy:

$$\Delta R^2 = R^2_{\text{weight}} - R^2_{\text{baseline}}$$

Success requires $\Delta R^2$ > 0.
